\chapter{多工位生产工装板设计实践}
\label{chap:case-programmer-extboard}

多工位生产工装板同时属于电源系统、信号接口、机械夹具和生产资料的一部分，不能按普通转接板处理。
本章说明如何把核心板接口、电源隔离、多路两线总线、工位机械尺寸、器件生命周期和生产输出统一到 KiCad 工程中。

\section{需求分解}

工装板需要直接连接控制核心板，并提供多个方向一致、可快速换片的操作工位。主要约束包括：
\begin{itemize}
  \item 各路 SCL/SDA 相互独立；
  \item 所有工位由全局使能控制，默认关闭；
  \item 插座侧电源可限流并自动放电；
  \item 未供电工位不能通过信号线反向供电；
  \item 工位布局、Pin~1 方向和丝印适合操作员重复装夹；
  \item 扩展板不能覆盖核心板较高的网络和 USB 连接器；
  \item 制造包包含原理图、装配、Gerber、钻孔、BOM、坐标和摘要。
\end{itemize}

\section{系统电路结构}

\begin{center}
\begin{tikzpicture}[node distance=6mm]
  \tikzset{hw/.style={draw=bookline,fill=bookpaper,rounded corners=1mm,
    minimum width=37mm,minimum height=10mm,align=center,font=\small}}
  \node[hw] (k4b) {控制核心板\\逻辑电源＋多路总线};
  \node[hw,right=of k4b] (power) {限流负载开关\\自动放电};
  \node[hw,below=of power] (switch) {双向模拟开关\\GLOBAL\_EN};
  \node[hw,right=of power] (supply) {Switched\_VCC\\本地去耦};
  \node[hw,right=of switch] (sockets) {四个烧录座\\独立 SCL/SDA};
  \draw[->,bookteal] (k4b)--node[above,font=\scriptsize]{VCC}(power);
  \draw[->,bookteal] (power)--(supply);
  \draw[->,bookteal] (k4b)--node[left,font=\scriptsize]{总线}(switch);
  \draw[->,bookteal] (switch)--(sockets);
  \draw[->,bookteal] (supply)--(sockets);
\end{tikzpicture}
\end{center}

电源开关限制所有工位的总电流并在关闭后释放插座电荷；双向模拟开关断开核心板与插座信号。
GLOBAL\_EN 使用下拉保证主控未初始化、复位或连接器松动时保持关闭。

电源与信号必须按同一安全状态设计，不能只关闭其中一路：

\begin{table}[htbp]
  \centering
  \caption{工装板电源与信号状态约束}
  \begin{tabular}{p{22mm}p{26mm}p{26mm}p{58mm}}
    \hline
    状态 & 工位电源 & 信号隔离 & 允许行为 \\
    \hline
    上电复位 & 关闭 & 断开 & 装卸器件、控制器初始化 \\
    待机 & 关闭 & 断开 & 检查夹具与操作条件 \\
    执行 & 开启且限流 & 接通 & 通信、写入和回读验证 \\
    故障 & 强制关闭 & 强制断开 & 记录错误，等待人工确认 \\
    \hline
  \end{tabular}
\end{table}

\section{电源开关设计}

负载开关需要满足输入电压、持续电流、限流精度、导通电阻、启动时间、反向电流和自动放电要求。
限流负载开关可通过外部电阻配置约束电流，并在输入、输出分别放置去耦电容。

限流值不是越大越好。过大的限流无法保护误插和短路，过小则可能在多工位同时启动或器件内部写入时误触发。
设计应依据单工位峰值、工位数量、连接器接触瞬态和温度降额确定：

\begin{equation}
  I_{limit,min} > N I_{socket,peak} + I_{margin}
\end{equation}

同时验证短路时器件结温、限流持续时间和 PCB 铜箔温升。

自动放电时间决定关闭使能后多久可以安全换片。样板必须测量插座电压下降曲线，而不能只依据典型放电电阻估算。

\section{信号隔离与上拉}

四组双向模拟开关分别连接核心板侧与插座侧 SCL/SDA，并由统一使能控制。
模拟开关选型需检查工作电压、导通电阻、带宽、关断漏电、掉电行为、钳位结构和封装 exposed pad 连接要求。

I2C 类开漏总线的上拉电阻由电容和目标上升时间决定：

\begin{equation}
  t_r \approx 0.8473 R_p C_b
\end{equation}

本地上拉接到可切换电源，使关闭工位时信号不会通过常电上拉反向供电。
如果某一路复用核心板已有上拉，必须确认其电压、阻值和关断状态，并在设计文档中标记这一外部依赖。

模拟开关的导通电阻会与上拉和器件输入形成附加压降，也会影响低电平裕量。
样板验证应同时观察 SCL/SDA 的上升时间、低电平、ACK 和插拔瞬态。

\section{去耦与回流}

每个烧录座具有本地去耦，负载开关和模拟开关也分别就近去耦。
双面地铜为主要回流路径，连接器地脚和器件地通过短路径与过孔连接。

小型数字板仍然需要关注回流连续性。SCL/SDA 不应跨越地铜开槽或狭窄颈部；连接器附近的地脚应为外部线缆和 ESD 电流提供明确返回路径。

\section{核心板连接器方向}

底面安装的板对板连接器容易产生顶视、底视和插合面三种方向混淆。
连接器定义应同时明确：
\begin{itemize}
  \item 核心板实物顶视的 Pin~1；
  \item 扩展板顶视焊盘 Pin~1；
  \item 从底视或插合面观察时的镜像关系；
  \item 位于底层丝印且紧邻真实焊盘的 Pin~1 标记。
\end{itemize}

任何连接器封装都应通过打印 1:1 纸样、实物叠放或机械模型验证。仅凭原理图左右顺序不能推断底面连接器的物理方向。

\section{烧录座封装与操作一致性}

烧录载板不是标准 DIP 封装，因此根据实物测量建立自定义 footprint，包括载板包络、列间距、同列针距、针截面和有效长度。
四个工位保持相同旋转方向，并在载板包络外永久显示 Pin~1 丝印，降低操作员反插概率。

自定义封装应同时表达：
\begin{itemize}
  \item 电气焊盘与钻孔；
  \item 实物包络和装配禁布区；
  \item 插拔方向与操作空间；
  \item 母排插入深度和夹持范围；
  \item Pin~1 及工位编号。
\end{itemize}

库封装与 PCB 内嵌实例必须使用相同局部坐标。更新库之前应比较焊盘位置，防止历史版本的符号翻转重新引入。

\section{机械舌片与支撑}

核心板的 RJ45 和叠层 USB 高于板对板母座，整板重叠会产生碰撞。
一种实用结构是采用短连接舌片，只在核心板边缘区域重叠，主要工位区域向外伸出。

三个安装孔与核心板下部对齐，外伸区域再使用独立铜柱支撑。
这种结构避免把插拔力和板重全部传给连接器，但必须验证铜柱高度、机箱落脚面和板面共面度。

从栅格图反推的机械坐标只能用于候选布局，量产前必须以 DXF、机械图或实物测量复核。
机械图中的“孔中心距板边”尺寸也不能误解为孔半径。

\section{器件生命周期与替代}

BOM 应明确标记关键器件的生命周期风险，并禁止使用名称相近但引脚或特性不同的变体。
这比仅记录阻值和封装更接近可量产 BOM。

关键器件条目应包含：
\begin{itemize}
  \item 精确制造商料号和封装代码；
  \item 引脚定义与 exposed pad 要求；
  \item 生命周期和采购状态；
  \item 允许替代料及必须复验的参数；
  \item 尚待实物确认的机械或电气条件。
\end{itemize}

\section{KiCad 源文件与生成资产}

以常见低速数字工装板为例，可以采用 1.6~mm FR4、双层铜和 ENIG 表面处理，但最终层叠和表面处理仍应由信号、电流、耐磨与成本要求决定。
Gerber、钻孔、装配坐标和评审 PDF 都是从源文件生成的派生产物。

工程曾经历多次候选版本，因此发布流程必须避免旧输出与当前源文件混合。
应以 KiCad 源文件和发布脚本为唯一来源，每次发布清空独立输出目录、重新生成全部文件，并对生产 ZIP 和内部文件生成摘要。

输出目录中的层命名、README 和 Gerber 层数也应自动与当前 stackup 核对，避免历史四层候选输出与当前双层源文件并存时造成误投板。

\section{ERC、DRC 与人工审查}

ERC 和 DRC 为零是发布必要条件，但不能发现所有错误。例如：
\begin{itemize}
  \item 底面连接器方向与实物镜像错误；
  \item 非标准载板针脚测量错误；
  \item 核心板高器件干涉；
  \item 负载开关实际采购变体错误；
  \item 关闭后的放电时间不满足换片安全；
  \item 总线波形因开关电阻和电容超出裕量。
\end{itemize}

因此候选发布还需要原理图评审、PCB 顶底评审、BOM 审查、1:1 机械检查和样板电气验证。

\section{样板验证矩阵}

\begin{enumerate}
  \item 检查 GLOBAL\_EN 上电默认关闭以及异常复位行为；
  \item 测量空载、单座和四座并发电流；
  \item 测量关闭后的 Switched\_VCC 放电时间；
  \item 观察四路 SCL/SDA 的高低电平、上升时间和 ACK；
  \item 验证空座、错位、热插拔误操作和短路保护；
  \item 执行高次数插拔，检查接触电阻和夹持力变化；
  \item 实物叠装确认连接舌、固定孔、RJ45、USB 和铜柱；
  \item 比较 BOM、坐标、丝印、Pin~1 和实际装配方向。
\end{enumerate}

\section{发布检查清单}

工装板发布包至少应完成以下闭环检查：
\begin{itemize}
  \item 原理图 ERC、PCB DRC 和网络表一致性均通过；
  \item 电源默认关闭、限流、放电和异常复位行为已经实测；
  \item 所有工位方向、Pin~1、通道编号和 UI 映射一致；
  \item 非标准封装经过 1:1 打印、实物叠放或机械模型复核；
  \item BOM 中的精确料号、替代料与生命周期状态已经审查；
  \item Gerber、钻孔、坐标、装配图和 BOM 均由同一版源文件重新生成；
  \item 生产 ZIP 内的层数、板厚、表面处理和版本标识与制造说明一致。
\end{itemize}

\section{工程方法总结}

专用工装设计中，电气、机械、操作和生产文件具有同等重要性。
应默认断电，使信号与电源同步隔离，固定工位方向，把实测机械信息写入 footprint，在 BOM 中记录生命周期风险，并把派生产物视为必须重新验证的发布结果。
